    \section{Resume}

    % skal være i kursiv
    % (PLACEHOLDER RESUMÉ, FÆRDIGGØR SENERE. TAK)
    Dette projekt har til hensigt at besvare, hvordan man kan hylde fejringen af Viborg Katedralskoles 100-års jubilæum for bygningen ved hjælp af et spil. Med udgangspunkt i den historiske begivenhed den 23. oktober 1944, hvor Gestapo gennemførte en razzia på skolen og arresterede 14 modstandsfolk, udvikles en digital oplevelse, der formidler skolens besættelseshistorie til nuværende og kommende elever i aldersgruppen 14-19 år. Projektet anvender Scrum og Unified Process som styringsmetoder og udvikles i Unreal Engine med fokus på multiplayer-elementer, der muliggør samarbejde mellem elever. Gennem co-design med lærere og elever fra Katedralskolen integreres musik, historiske objekter og dilemma-baserede fortællinger. Resultatet er en spil-prototype, der demonstrerer, hvordan en historisk bygning kan fremvises i et historisk perspektiv, hvordan elever kan engagere sig gennem kreative tilføjelser, og hvordan en digital oplevelse kan designes med udgangspunkt i en konkret historisk begivenhed. Projektet bidrager med empirisk indsigt i anvendelsen af serious games til historieformidling i en skolekontekst.
